\subsection{The x- and y-intercepts of a Quadratic Function} 
It can be helpful, in planning a graph, to find its $x$- and $y$-intercepts. Recall the following definitions.
\begin{definition}[Intercepts]
The \term{$x$-intercept(s)} of a graph is(are) the point(s) at which the graph intersects the $x$-axis. The \term{$y$-intercept(s)} of a graph is(are) the point(s) at which the graph intersects the $y-$axis.
\end{definition}
The graph of a function only has one \makeblanksmall-intercept, but might have more than one \makeblanksmall-intercept.
\begin{itemize}
\item To find the $y$-intercept, plug in $\makeblanksmall=0$ and \makeblank[1in].
\item To find the $x$-intercept, set $\makeblanksmall=0$ and \makeblank[1in].
\end{itemize}
\begin{Example}
Find the $x$- and $y-$ intercepts of the graph of $g(x)=x^2-x-6$.
\begin{lpic}{}{5}
\foreach \x/\y/\txt in {0/1/{Set $x=0$},0/3/{$y$-intercept:},1/1/{Set $g(x)=0$},1/5/{$x$-intercepts:}}
	\regtextleft{\x*\value{cols}+0.25}{-\y}{\txt};
\end{lpic}
We can use this information to help us sketch the graph. Note that the $y$-intercepts have the same $y$-coordinate, so the axis of symmetry is halfway between them.
\begin{lpic}{}{11}
\regtextleft{0.25}{-1}{Axis of Symmetry:};
\regtextleft{0.25}{-2}{Vertex:};
\regtext{1}{-3}{$x$};
\regtextleft{3}{-3}{$y=g(x)$};
\draw (0,-3) -- (6,-3);
\draw (2,-2) -- (2,-11);
\end{lpic}
\end{Example}